For fixed $S$ and $T$, the states of the system are classified by total angular momentum $J$ and parity.
As we know, $T=0$ corresponds to states with symmetric wave functions $\psi$, whereas for $T=1$ these functions are antisymmetric.
On the other hand, $S$ fixes the symmetry of the wave function in the spin variables: it is symmetric for $S=1$ and antisymmetric for $S=0$.
Consequently, the pair $S,T$ also fixes the symmetry in the spatial variables, and hence the parity of the state.
Thus, a system with isotopic spin $T=0$ can have only even triplet states ($S=1$) or odd singlet states ($S=0$); for $T=1$, the states are odd triplets or even singlets.

Since the spin vector is not conserved, orbital angular momentum is not generally conserved either; only the sum $\mathbf J=\mathbf L+\mathbf S$ is conserved.
Its magnitude $L$ may nevertheless be conserved simply because fixed values of $J$, $S$, and parity (or of $J$, $S$, and $T$) can be compatible with just one definite value of $L$; recall that the parity of a two-particle system is $(-1)^L$.
